\section{性能分析}
\label{sec:performance_analysis}

该算法的优缺点需要通过仿真和安全测试验证。因此,本部分介绍算法的仿真及性能测试结果。仿真实验在AMD Ryzen 7 4800H CPU,16 GB 内存,Windows 11的，MATLAB 2024b环境下进行，对多张尺寸不同的图片进行了仿真测试。

\subsection{仿真结果}
从USC-SIPI图像数据库中选择了3张$256\times256$、一张$512\times512$和一张$128\times128$的图像进行仿真测试。图[fig]展示了采用所提加密算法得到的加密图像和解密图像。可以看出，加密图像为不规则噪声信号图像，难以从中获取原始图像的任何信息；解密图像则与原始图像一致。

\subsection{秘钥空间分析}
所有可能秘钥的总数称为秘钥空间。秘钥空间越大，算法的安全性越高。Alvarez[ref]指出,安全的加密算法的秘钥空间至少应为$2^{100}$，且加密算法的秘钥空间越大，其抵抗暴力破解的能力就越强。本文所提加密算法可以是384位哈希值H，或是五个参数$\{\mu,v,\alpha,\beta,A_0\}$。(1)若秘钥为384位哈希值H，则秘钥空间为$2^{384}$；(2)若秘钥为五个参数$\{\mu,v,\alpha,\beta,A_0\}$，在计算机精度为$10^{-15}$的情况下，则秘钥空间为$10^{5 \times 15} = 10^{75} > 2^{249}$。考虑到攻击者可能通过暴力破解或其他方式获取秘钥，最终秘钥空间应为两者的最小值，即$10^{75} > 2^{249}$，因此本文所提加密算法的秘钥空间足够大，能够抵抗暴力破解攻击。

\subsection{像素分布分析}
一张像素分布均匀的图像可有效抵抗统计分析攻击，通过直方图能够直观的反映出图片的像素分布情况。一幅视觉上有意义的图像通常呈现出不均匀的像素分布，其直方图有明显的峰谷特征。而一幅加密后的图像应呈现出均匀的像素分布，其直方图应接近于均匀分布。图[fig]展示了原始图像和加密图像的直方图对比。可以看出，加密图像的直方图呈现出较为均匀的分布，各区间高度一致。此外我们引入了直方图方差(HV)，定义式见~\eqref{eq:hvar},以及$\chi^2$,定义式见~\eqref{eq:chi2},用以进一步分析加密图像的像素分布情况。
\begin{equation}
\label{eq:hvar}
\mathrm{Hvar}(X) = \frac{1}{n^2} \sum_{i=1}^{n} \sum_{j=1}^{n} \frac{(z_i - z_j)^2}{2}
\end{equation}
其中，$z_i$ 或 $z_j$ 表示像素值 $i$ 或 $j$ 的统计频数。

\begin{equation}
\label{eq:chi2}
\chi^2 = \sum_{i=0}^{255} \frac{(f_i - f_0)^2}{f_0}
\end{equation}
其中，$f_0 = (M \times N) / 256$ 表示理想均匀分布下像素值的期望出现频率（$M \times N$ 为图像尺寸），$f_i$ 为灰度值 $i$ 的实际出现频数。当显著性水平 $\alpha$ 设为 $0.005$（自由度为 $255$）时，卡方分布的临界阈值为 $293.24783$。

与直方图相比，$HV$ 能更准确直观地反映像素分布的均匀性，表~\ref{tab:hv_comparison}展示了我们算法与其他算法加密前后像素分布的对比；若密文图像计算所得的 $\chi^2$ 均小于 $293.24783$，则说明密文图像通过了显著性检验，像素值呈现出高度均匀的分布。

\begin{table}[htbp]
\centering
\caption{The results of $HV$ test and the comparison with the related algorithms.}
\label{tab:hv_comparison}
\begin{tabular}{llcc}
\toprule
Algorithms & Images & Plain image & Encrypted image \\
\midrule
proposed               & img1.png(Sailboat) & 895231.41  & 974.35 \\
                       & img2.png(Baboon) & 407453.85  & 997.08  \\
                       & img3.png(Lena) & 950136.46  & 932.62 \\
                       & img4.png(Peppers) & 1363997.77 & 976.16  \\
                       & img5.png(Aerial) & 727258.66  & 964.43 \\
\textbf{Avg. proposed} & --       & --         & 968.93 \\
Ref1                   & --       & --         & 1014.33 \\
Ref2                   & --       & --         & 1013.00 \\
Ref3                   & --       & --         & 1004.02  \\
Ref4                   & --       & --         & 1041.58 \\
\bottomrule
\end{tabular}
\end{table}